\documentclass[11pt,letterpaper]{article}
\usepackage[margin=2.5cm]{geometry}
\usepackage[utf8]{inputenc}
\usepackage[T1]{fontenc}
\usepackage{lmodern}
\usepackage{amsmath,amssymb,amsthm,mathtools,bm}
\usepackage{enumitem}
\usepackage{booktabs}
\usepackage{tikz}
\usetikzlibrary{arrows.meta,patterns,positioning,calc}
\usepackage{pgfplots}
\pgfplotsset{compat=1.16}
\usepackage{titlesec}
\usepackage[hidelinks]{hyperref}

\setlength{\parindent}{0pt}
\setlength{\parskip}{0.5em}
\titleformat{\section}{\large\bfseries}{\thesection.}{0.6em}{}
\titleformat{\subsection}{\normalsize\bfseries}{\thesubsection.}{0.6em}{}

\theoremstyle{definition}
\newtheorem{definition}{Definition}[section]
\newtheorem{example}{Example}[section]
\newtheorem{exercise}{Exercise}[section]
\newtheorem{theorem}{Theorem}[section]
\newtheorem{remark}{Remark}[section]

\newcommand{\R}{\mathbb{R}}
\newcommand{\N}{\mathbb{N}}
\newcommand{\Z}{\mathbb{Z}}
\newcommand{\U}{\mathcal{U}}
\newcommand{\E}{\mathbb{E}}
\newcommand{\Var}{\operatorname{Var}}
\newcommand{\Cov}{\operatorname{Cov}}
\newcommand{\Normal}{\mathcal{N}}
\newcommand{\dd}{\,\mathrm{d}}
\newcommand{\T}{^{\top}}
\renewcommand{\P}{P}
\newcommand{\set}[1]{\left\{#1\right\}}
\newcommand{\given}{\mid}
\newcommand{\note}[1]{\textit{#1}}

% --- Homework exercises (tinted box) and highlighted derivation steps ---
\usepackage{tcolorbox}
\tcbuselibrary{skins,breakable}
\definecolor{hwfill}{RGB}{235,243,252}
\definecolor{hwframe}{RGB}{31,78,140}
\definecolor{newterm}{RGB}{196,78,0}
\newtcolorbox{homework}[1]{enhanced, breakable,
  colback=hwfill, colframe=hwframe, coltitle=white,
  fonttitle=\bfseries, boxrule=0.6pt, arc=2pt,
  left=6pt, right=6pt, top=4pt, bottom=4pt,
  before skip=10pt, after skip=10pt,
  title={Exercise (Homework #1)}}
\newcommand{\soln}{\par\smallskip\noindent\textbf{Solution.}\ }
% \new{...} marks what is new in a derivation step: a substituted or simplified term.
\newcommand{\new}[1]{{\color{newterm}#1}}

\title{Lecture 8: Markov Chain Monte Carlo}
\author{}\date{}

\begin{document}
\maketitle

\section{Motivation: sampling instead of integrating}
For non-conjugate families in the Bayes rule
\[
p(\theta \mid y) = \frac{p(y \mid \theta)\,p(\theta)}{p(y)},
\]
the main issue is to compute
\[
p(y) = \int \mathrm{d}\theta'\,p(y \mid \theta')\,p(\theta').
\]
Also, it is hard to make predictions using
\[
p(\tilde y \mid y) = \int \mathrm{d}\theta\,p(\tilde y \mid \theta)\,p(\theta \mid y).
\]
Instead of finding the integral, we can sample $N$ $\theta$s from $p(\theta \mid y)$. This is, we can generate $\set{\theta^{(i)}}_{i=1}^{N}$, where
\[
\theta^{(i)} \sim p(\theta \mid y),
\]
and then use
\[
p(\tilde y \mid y) = \E_{\theta \sim p(\theta \mid y)}\left[p(\tilde y \mid \theta)\right] \approx \frac1N\sum_{i=1}^{N}p(\tilde y \mid \theta^{(i)}).
\]
But notice that I also need $\tilde y$. So, what is actually done is to sample an
\[
\tilde y^{(i)} \sim p(\tilde y \mid \theta^{(i)})
\]
for each $\theta^{(i)} \sim p(\theta \mid y)$, where $y$ is known. Then we have a set $\set{\tilde y^{(i)}}_{i=1}^{N}$.

We never get $p(\tilde y \mid y)$ in closed form; we only get samples from it, from which any expectation can be approximated.

\begin{example}[Sampling from the exponential distribution]
The exponential PDF is
\[
f(x) = \lambda e^{-\lambda x}, \qquad x > 0,\ \lambda > 0.
\]
We want to get samples from $f(x)$. This is, $x \sim f(x)$. The CDF is
\[
F_X(x) = P(X \le x) = \int_0^x \new{\lambda e^{-\lambda u}}\dd u = \new{1 - e^{-\lambda x}},
\]
where $0 \le F_X(x) \le 1$. For sampling, we can generate
\[
u \sim \operatorname{Uniform}(0,1)
\]
and make $F_X(x) = u$. Thus $e^{-\lambda x} = 1 - u$, and then we have
\[
-\lambda x = \new{\log}(1 - u) \ \Rightarrow\ x = \new{-\frac{1}{\lambda}}\log(1 - u).
\]
\end{example}

\begin{example}[Sampling from a multivariate normal]
To sample from $\Normal(x \mid \mu, \Sigma)$ we can use
\[
x = Lz + \mu, \qquad z \sim \Normal(0, I),
\]
where $LL^{\top} = \Sigma$ is the Cholesky factorization of $\Sigma$. Notice that
\[
\E(x) = L\,\new{\E(z)} + \new{\E(\mu)} = \new{\mu}
\]
and
\[
\Cov(x) = \E\left((x - \mu)(x - \mu)^{\top}\right) = \E\left[\new{Lzz^{\top}L^{\top}}\right] = \new{L}\underbrace{\E(zz^{\top})}_{\new{I}}\new{L^{\top}} = \new{\Sigma},
\]
and therefore $\Cov(x) = LL^{\top} = \Sigma$.

Sampling $z \sim \Normal(0, I)$ is done using, for example, the Box--Muller method in $O(d)$, where $d$ is the size of the vector $x$. The overall cost of sampling $x \sim \Normal(x \mid \mu, \Sigma)$ is $O(d^3)$, which includes the cost of computing the Cholesky factorization ($O(d^3)$).
\end{example}

%=====================================================
\section{Metropolis--Hastings}
It is a Markov chain Monte Carlo (MCMC) method used for obtaining a sequence of random samples from a probability distribution from which direct sampling is difficult. MCMC is a class of algorithms used to draw samples from a probability distribution.

New samples are added to the sequence in two steps:
\begin{enumerate}[label=(\roman*)]
\item a new sample is proposed based on the previous one (this is the Markovian part), and
\item the proposed sample is either added to the sequence or rejected depending on the value of the probability distribution at that point.
\end{enumerate}

\subsection{Formal derivation}
The problem is that we need to sample from a target distribution
\[
\pi(x) = \frac{\tilde\pi(x)}{Z},
\]
where we can only evaluate $\tilde\pi(x)$ and $Z$ is unknown.

To generate samples using the Metropolis--Hastings algorithm, we use a \emph{proposal distribution} $r(y \mid x)$ that we choose and from which it is easy to sample.

Let $x$ be a current state, so $X_t = x$. We draw a candidate $y$ from a proposal distribution $r(y \mid x)$. If we accept it, then $X_{t+1} = y$. Otherwise, $X_{t+1} = x$.

How do we decide whether to accept the candidate $y$?
\begin{enumerate}
\item Initialize $x \in \R^n$.
\item Sample $y \sim r(\cdot \mid x)$.
\item Sample $u \sim \operatorname{Uniform}(0,1)$.
\item If $u \le \alpha(y \mid x)$, accept $y$ and set $x \leftarrow y$. Otherwise, keep $x$.
\item Return to step 2.
\end{enumerate}
where
\[
\alpha(y \mid x) = \min\Bigg\{1,\ \underbrace{\frac{\tilde\pi(y)\,r(x \mid y)}{\tilde\pi(x)\,r(y \mid x)}}_{\rho}\Bigg\}.
\]
If $\rho \ge 1$, we always accept $y$. If $\rho < 1$, we accept $y$ with probability $\rho$.

\subsection{Why $\alpha(y \mid x)$?}
Suppose that at step $t$ we have $X_t \sim \pi(\cdot)$. We would like the next state to have the same distribution, so that $X_{t+1} \sim \pi(\cdot)$ as well. Can we find a transition probability $p(y \mid x)$ that accomplishes that? This is,
\[
P(X_{t+1} = y) = \sum_x p(y \mid x)\,\pi(x).
\]
We want a transition rule such that
\[
\pi(y) = \sum_x p(y \mid x)\,\pi(x).
\]
This is, given many copies of the chain, each state distributed according to $\pi(\cdot)$: after every copy takes one step, what fraction do we expect to find at $y$?

One way of having stationarity is \emph{detailed balance}:
\[
\pi(x)\,p(y \mid x) = \pi(y)\,p(x \mid y).
\]
Thus
\[
\sum_x p(y \mid x)\,\pi(x) = \sum_x \new{\pi(y)\,p(x \mid y)} = \new{\pi(y)}\underbrace{\sum_x p(x \mid y)}_{\new{1}} = \pi(y).
\]

The probability of a candidate $Y = y$ being accepted ($A$), given $X_t = x$, is
\[
P(Y = y, A \mid X_t = x) = P(A \mid Y = y, X_t = x)\,P(Y = y \mid X_t = x).
\]
That is the same as saying
\[
\underbrace{P(X_{t+1} = y \mid X_t = x)}_{\new{p(y \mid x)}} = \underbrace{P(A \mid Y = y, X_t = x)}_{\new{\alpha(y \mid x)}}\,\underbrace{P(Y = y \mid X_t = x)}_{\new{r(y \mid x)}}.
\]
Finally, for $y \neq x$,
\[
p(y \mid x) = \alpha(y \mid x)\,r(y \mid x).
\]
(For $y = x$, rejections also add probability, but detailed balance holds trivially when $y = x$.)
Now, I replace $p(y \mid x)$ into the detailed balance equation
\[
\pi(x)\,p(y \mid x) = \pi(y)\,p(x \mid y),
\]
\[
\pi(x)\,\new{\alpha(y \mid x)\,r(y \mid x)} = \pi(y)\,\new{\alpha(x \mid y)\,r(x \mid y)},
\]
where $x \to y$ means the chain starts at $x$ and proposes $y$, and
\[
\alpha(x \mid y) = P(A \mid Y = x, X_t = y) \qquad (y \to x:\ \text{starts at } y \text{ and proposes } x).
\]
Then
\[
\alpha(y \mid x) = \underbrace{\frac{\pi(y)\,r(x \mid y)}{\pi(x)\,r(y \mid x)}}_{\new{\rho}}\,\alpha(x \mid y)
\ \Rightarrow\
\alpha(y \mid x) = \new{\rho}\,\alpha(x \mid y).
\]
Since $\pi(y)/\pi(x) = \tilde\pi(y)/\tilde\pi(x)$, $Z$ is not needed. MH makes the larger of the two acceptance probabilities equal to $1$: if $\rho < 1$, it takes $\alpha(x \mid y) = 1$, so $\alpha(y \mid x) = \rho$; if $\rho \ge 1$, it takes $\alpha(y \mid x) = 1$ (and then $\alpha(x \mid y) = 1/\rho$). Finally,
\[
\alpha(y \mid x) = \min(1, \rho).
\]

%=====================================================
\section{Hamiltonian Monte Carlo (HMC)}
HMC is particularly suited for continuous differentiable density distributions.

HMC introduces momentum variables $q$ and draws from a joint distribution
\[
p(q, \theta) = p(q \mid \theta)\,p(\theta).
\]
For most applications,
\[
p(q \mid \theta) = p(q) = \Normal(0, \Sigma).
\]
Normally $\Sigma = I$ or it is estimated from warm-up.

The negative log of $p(q, \theta)$ defines a Hamiltonian:
\[
H(q, \theta) = -\log\big(p(q, \theta)\big) = \new{-\log\big(p(q \mid \theta)\big) - \log\big(p(\theta)\big)},
\]
\[
H(q, \theta) = \new{K(q \mid \theta)} + \new{V(\theta)},
\]
where $K(q \mid \theta) = -\log[p(q \mid \theta)]$ is called the \emph{kinetic energy} and $V(\theta) = -\log(p(\theta))$ is called the \emph{potential energy}.

\subsection{Generating transitions}
Starting from the current value of $\theta$, a transition to a new state is generated in two stages:
\begin{enumerate}
\item A value for $q$ is drawn, independently of the current value of $\theta$:
\[
q \sim p(q) = \Normal(q \mid 0, \Sigma).
\]
\item $q$ and $\theta$ are evolved according to
\[
\frac{\mathrm{d}\theta}{\mathrm{d}t} = \nabla_q H = \nabla_q \new{K(q)},
\qquad
\frac{\mathrm{d}q}{\mathrm{d}t} = -\nabla_\theta H = -\nabla_\theta \new{V(\theta)}.
\]
Thus
\[
\frac{\mathrm{d}\theta}{\mathrm{d}t} = \nabla_q K(q), \qquad \frac{\mathrm{d}q}{\mathrm{d}t} = -\nabla_\theta V(\theta).
\]
\end{enumerate}

\subsection{Leapfrog integrator}
This algorithm takes discrete steps of some small time interval $\varepsilon$. Leapfrog begins by drawing a fresh momentum value
\[
q \sim \Normal(q \mid 0, \Sigma)
\]
and then updates are done as
\begin{align*}
q &\leftarrow q - \frac{\varepsilon}{2}\nabla_\theta V(\theta),\\
\theta &\leftarrow \theta + \varepsilon\,\nabla_q K(q),\\
q &\leftarrow q - \frac{\varepsilon}{2}\nabla_\theta V(\theta).
\end{align*}
By applying $L$ leapfrog steps, a total time $L\varepsilon$ is simulated. The resulting state after $L$ repetitions is denoted $(\tilde q, \tilde\theta)$.

For the Gaussian momentum, $K(q) = \frac12 q^{\top}\Sigma^{-1}q$ up to a constant, so the position step is $\theta \leftarrow \theta + \varepsilon\,\Sigma^{-1}q$: the position moves \emph{along} the momentum, with a plus sign.

\subsection{Metropolis accepting step}
The leapfrog integrator does not conserve $H$ exactly, so the proposal $(\tilde q, \tilde\theta)$ is accepted with probability
\[
\alpha(q, \theta, \tilde q, \tilde\theta) = \min\Big[1,\ \exp\big(\new{H(q, \theta) - H(\tilde q, \tilde\theta)}\big)\Big].
\]
This is the Metropolis ratio for the joint target $p(q, \theta) \propto e^{-H(q, \theta)}$:
\[
\frac{p(\tilde q, \tilde\theta)}{p(q, \theta)} = \frac{e^{-H(\tilde q, \tilde\theta)}}{e^{-H(q, \theta)}} = \new{e^{H(q, \theta) - H(\tilde q, \tilde\theta)}}.
\]
If $\tilde\theta$ is accepted, it is the next state; otherwise the chain stays at $\theta$. The momentum is discarded either way, and a fresh one is drawn at the next iteration. If the dynamics were simulated exactly, $H$ would be conserved and every proposal would be accepted.

\subsection{Aside: where Hamilton's equations come from}
In classical mechanics a system with coordinates $q$ is described by a Lagrangian $L(q, \dot q, t)$, and its path makes the action
\[
S = \int L(q, \dot q, t)\dd t
\]
stationary. This gives the \emph{Euler--Lagrange equations}
\[
\frac{\partial L}{\partial q} - \frac{\mathrm{d}}{\mathrm{d}t}\frac{\partial L}{\partial \dot q} = 0.
\]
Hamilton then defined, with momenta $p_i = \partial L/\partial \dot q_i$,
\[
H(q, p, t) = \sum_i \dot q_i\,p_i - L(q, \dot q, t),
\]
and the Euler--Lagrange equations become the first-order system
\[
\frac{\mathrm{d}q}{\mathrm{d}t} = \frac{\partial H}{\partial p}, \qquad \frac{\mathrm{d}p}{\mathrm{d}t} = -\frac{\partial H}{\partial q}.
\]
\emph{Notation:} in this aside $q$ is the position and $p$ the momentum, as in physics. In HMC above the position is $\theta$ and the momentum is $q$, so $\mathrm{d}\theta/\mathrm{d}t = \nabla_q H$ and $\mathrm{d}q/\mathrm{d}t = -\nabla_\theta H$ are the same equations.

%=====================================================
\section{Convergence: the Gelman--Rubin statistic}
Run $M$ chains with $N$ draws each (after warm-up), from dispersed starting points. For a scalar quantity $\theta$, let
\begin{itemize}[nosep]
\item $\bar\theta_m$ be the average of chain $m$,
\item $s_m^2$ be the variance of chain $m$,
\item $\bar\theta = \dfrac1M\displaystyle\sum_{m=1}^{M}\bar\theta_m$ be the average over all chains.
\end{itemize}
The \emph{within-chain} variance $W$ and the \emph{between-chain} variance $B$ are
\[
W = \frac1M\sum_{m=1}^{M}s_m^2, \qquad
B = \frac{N}{M-1}\sum_{m=1}^{M}\big(\bar\theta_m - \bar\theta\big)^2,
\]
where $B$ measures how much the chains differ from each other. They are combined into an estimate of the posterior variance,
\[
\hat V = \frac{N-1}{N}\,W + \frac1N\,B,
\]
and the \emph{potential scale reduction factor} is
\[
\hat R = \sqrt{\frac{\hat V}{W}}.
\]
If the chains have mixed, they all explore the same distribution, $B$ is comparable to $W$, and $\hat R \approx 1$: this is the sign of convergence. A value clearly above $1$ means the chains have not yet agreed, and they should be run longer.

%=====================================================
\section*{References}
\begin{itemize}
\item Gelman, A., Carlin, J.~B., Stern, H.~S., Dunson, D.~B., Vehtari, A., and Rubin, D.~B. (2013). \emph{Bayesian Data Analysis} (3rd ed.). CRC Press. Chapters 10--12: simulation, Metropolis--Hastings and Hamiltonian Monte Carlo; Section 11.4: assessing convergence with $\hat R$; Section 12.4: Hamiltonian Monte Carlo and the leapfrog steps.
\item Chib, S., and Greenberg, E. (1995). Understanding the Metropolis--Hastings algorithm. \emph{The American Statistician}, 49(4), 327--335. The detailed-balance derivation of the acceptance probability.
\item Neal, R.~M. (2011). MCMC using Hamiltonian dynamics. In S.~Brooks, A.~Gelman, G.~Jones and X.-L.~Meng (Eds.), \emph{Handbook of Markov Chain Monte Carlo}. CRC Press. Hamiltonian dynamics and the leapfrog integrator.
\item Betancourt, M. (2017). A conceptual introduction to Hamiltonian Monte Carlo. \url{https://arxiv.org/abs/1701.02434} Intuition for HMC.
\item Kuleshov, V., and Ermon, S. \emph{CS228 Notes: Probabilistic Graphical Models}. \url{https://ermongroup.github.io/cs228-notes/} Section on sampling methods.
\end{itemize}

\end{document}
